
\documentclass[a4paper]{article}
\usepackage[french]{babel}
\usepackage{geometry}

\usepackage{tikz}
\usepackage{pifont}
\usepackage{graphicx}

\newlength\cardheight
\newlength\cardwidth
\setlength\cardheight{87mm}
\setlength\cardwidth{56mm}

\newlength\cardhspace
\newlength\cardvspace
\setlength\cardhspace{56.5mm}

\def\tabularxcolumn#1{m{#1}}
\pgfmathsetmacro{\cardroundingradius}{4mm}
\pgfmathsetmacro{\striproundingradius}{3mm}


\begin{document}

\newgeometry{margin=10mm}
\thispagestyle{empty}

\newcommand{\stripcolor}{cyan}
\pgfmathsetmacro{\stripwidth}{1.2}
\pgfmathsetmacro{\strippadding}{0.1}

\pgfmathsetmacro{\textpadding}{0.3}

\newcommand{\topcaption}{Feu vert}
\newcommand{\bottomcaption}{Feu rouge}
\pgfmathsetmacro{\ruleheight}{0.1}

\newcommand{\stripfontsize}{\Huge}
\newcommand{\captionfontsize}{\LARGE}
\newcommand{\textfontsize}{\normalsize }

\clearpage

\pgfmathsetmacro{\cardwidth}{5}
\pgfmathsetmacro{\cardheight}{8}





\newcommand{\gamecard}[3]{%
 \begin{tikzpicture}
     \draw[rounded corners=\cardroundingradius] (0,0) rectangle (\cardwidth,\cardheight);
     \fill[\stripcolor,rounded corners=\striproundingradius] (\strippadding,\strippadding) rectangle (\strippadding+\stripwidth,\cardheight-\strippadding) node[rotate=90,above left,black,font=\stripfontsize] {#1};
     \node[text width=(\cardwidth-\strippadding-\stripwidth-2*\textpadding)*1cm,below right,inner sep=0] at (\strippadding+\stripwidth+\textpadding,\cardheight-\textpadding) 
     {   {\captionfontsize \topcaption}\\ 
         {\textfontsize \textit{#2}}\\
         \tikz{\fill (0,0) rectangle (\cardwidth-\strippadding-\stripwidth-2*\textpadding,\ruleheight);}\\
         {\captionfontsize \bottomcaption}\\ 
         {\textfontsize \textit{#3}}\\
     };
 \end{tikzpicture}
 

}

% carte initiale
 \begin{tikzpicture}
     \draw[rounded corners=\cardroundingradius] (0,0) rectangle (\cardwidth,\cardheight);
        \node[text width=(\cardwidth-\strippadding-2*\textpadding)*1cm,below right,inner sep=0] at (\strippadding+\textpadding,\cardheight-\textpadding) 
    { \centering {\Huge  {Healthcheck card deck}}
    {\paragraph{traduction: Ajiro.fr}
    Juillet 2016}
    
	{\paragraph{Credits:}   
    Health check model: Henrik Kniberg \& Kristian Lindwall,
    with help from the other agile coaches at Spotify
    
    Feel free to spread/modify/reuse this model!Creative Commons Attribution-ShareAlike 
    }
    
    };
       
 \end{tikzpicture}

\clearpage

% carte Ajiro
 \begin{tikzpicture}
     \draw[rounded corners=\cardroundingradius] (0,0) rectangle (\cardwidth,\cardheight);
        \node
        [text width=(0,below right,inner sep=0] at (0.5,\cardwidth+1.5)
    {  
    \includegraphics[scale =0.4, angle =90 ]{logo-ajiro.jpg};};
 \end{tikzpicture}

\clearpage

%
\gamecard{Livraison de valeur}
  {Nous livrons des produits extraordinaires! Nous en sommes fiers et nos clients en sont particulièrement contents.}
  {Ce que nous livrons est nul. Nous en avons honte. Nos clients nous d\'etestent.}
%


\clearpage
\gamecard{Facilit\'e de livraison}
  {Le processus de livraison est simple, sûr, indolore et quasiement totalement automatis\'e.}
  {Le processus de livraison est risqu\'e, douloureux, n\'ecessite beaucoup d'interventions manuelles, qui prennent \'enorm\'em\'enet de temps.}
%
\clearpage
\gamecard{Fun}
  {Nous aimons aller travailler le matin, et avons beaucoup de fun à travailler ensemble.}
  {Ennui\ldots{}.}
%
\clearpage

\gamecard{Bonne sant\'e du code}
  {Nous sommes fiers de la qualit\'e de notre code. Il est propre, facile à lire et a une grande couverture de tests.}
  {Notre code est un tas de fumier et notre dette technique croît de fa\c{c}on incontrôlable.}
  \clearpage
\gamecard{Apprentissage}
  {Nous apprenons à chaque instant quelque chose d'int\'eressant.}
  {Nous n'avons jamais le temps d'apprendre quoi que ce soit.}
  
  \clearpage
\gamecard{Mission}
  {Nous savons exactement pourquoi nons sommes là, et nous sommes passionn\'es en connaissance de cause.}
  {Nous ne savons pas pouquoi nous sommes là, il n'y a pas de sch\'ema g\'en\'eral  ou d'orientation globale. Nos ``missions'' sont totalement floues et d\'emotivantes.}
%

\clearpage
%
\gamecard{Acteurs\ldots{} ou pions\ldots{}}
 {Nous contrôlons notre destin. Nous d\'ecidons ce que nous d\'eveloppons et comment nous le d\'eveloppons.}
 {Nous ne sommes que des pions sur un \'echiquier, et nous n'avons aucune influence sur ce que nous construisons et comment nous le d\'eveloppons.}
 
 
 \clearpage
\gamecard{V\'elocit\'e}
  {Nous faisons ce que nous avons d\'ecid\'e, sans délai, sans attendre Pas d'attente, pas de d\'elai.}
  {Nous ne finissons jamais rien. Nous sommes syst\'ematiquement interrompus et bloqu\'es par les d\'ependances critiques d'autres produits.}
  \clearpage
\gamecard{Processus adapt\'e}
  {Notre fa\c{c}on de travailler nous correspond tout à fait !}
  {Notre fa\c{c}on de travailler est ridicule.}
%
\clearpage
\gamecard{Support}
  {Nous obtenons toujours du support et de l'aide de nos sponsors dès que nous le demandons.}
  {Nous sommes r\'egulièrement bloqu\'es car nous ne pouvons obtenir du soutien et de l'aide de la part de nos sponsor quand nous le demandons.}
%

\clearpage

\gamecard{Travail d'\'equipe}
  {Nous sommes un groupe qui a vraiment pris une forme d'\'equipe, avec un niveau de collaboration extraodinaire.}
  {Nous sommes une bande d'individus qui ne se connaissent pas et nous ne sous soucions pas de savoir ce que les autres font.}
%



\restoregeometry

\end{document}



